%% sec7_discussion.tex -- Section 7: Discussion and Conclusion

\section{Discussion}
\label{sec:discussion}

The classical literature on lotto designs \cite{Furedi1996, Li2002,
Gordon1995} seeks the minimum number of tickets required to guarantee a
prize.  Our probabilistic framework addresses a complementary question:
given exactly $n$ tickets (typically far fewer than the deterministic
minimum), how should they be arranged to maximise the probability of
winning?  The two objectives align when $n \ge C(v,k,t)$, but for practical
budgets they lead to different optimisation problems.  Theorem~\ref{thm:two-ticket-monotone}
and Proposition~\ref{prop:mutual-excl} show that the probabilistic problem
has a cleaner structure in the high-$k$ regime: at near-minimum overlap the
two winning events are mutually exclusive, so $P_{\mathrm{win}} = 2p$
exactly.  No analogous result appears in the deterministic covering literature.

\citet{Liu2024} study a related probabilistic formulation for the UK 6/59
lottery and find that balanced pairwise overlap, rather than minimum overlap,
maximises expected return across all prize tiers.  Their finding is not in
contradiction with ours.  For the top prize level only, which is what we
address, minimum overlap is optimal for $n=2$ as proved in
Theorem~\ref{thm:two-ticket-monotone}.  For a multi-tier prize structure,
the optimal strategy may involve non-minimum overlap at certain tiers; we
leave this extension for future work.

The framework applies to any $(v,k)$-lottery with a single prize threshold.
Three natural extensions illustrate the range of parameters.  In Lotto 6/49
(Canada, Germany), where $v=49$, $k=6$, $m^*=0$, tickets can be disjoint
and the mutual-exclusivity result follows trivially.  The same holds for the
Brazilian Mega-Sena ($v=60$, $k=6$) and Quina ($v=80$, $k=5$).  The
Lotof\'{a}cil is distinguished by its ratio $k/v = 15/25 = 0.6$, which
forces non-trivial overlap and makes the optimisation problem genuinely hard.

\begin{remark}[Expected value]
  In a model with a fixed prize $V$ and ticket cost $c$, the expected return
  of an $n$-ticket collection is $V \cdot P_{\mathrm{win}}(\mathcal{A}) - nc$.
  Because $V$ and $nc$ are constants, maximising $P_{\mathrm{win}}$ is
  equivalent to maximising expected return: our method reduces the expected
  loss for a given budget, even though it cannot make expected return positive.
  The Lotof\'{a}cil's actual return rate is approximately 45\% of revenue
  \citep{CEF2024}, so the expected return on a single ticket of cost~$c$ is
  roughly $-0.55c$, regardless of the numbers chosen.  In the real game,
  prizes at each tier are pooled and divided among all winners, so the
  expected return also depends on how many other players hold the same
  winning combination.  Our method addresses only the $P_{\mathrm{win}}$
  component; the prize-division component would require data on the
  distribution of bets across all players, which we do not use.
\end{remark}

There are two main limitations to the current analysis.  First, the model
ignores prize division.  In practice, a ticket that avoids popular number
combinations earns a higher expected payout even with the same prize
probability, because fewer other players hold the same ticket.  Incorporating
this would require data on the distribution of all bets placed, which we
do not use.  Second, Conjecture~\ref{conj:general} remains unproved for
$n \ge 3$.  A complete proof would require monotone-coupling results for
multivariate hypergeometric distributions that, to our knowledge, are not
available in the literature.  Establishing such results would also have
implications for other statistical problems where multiple dependent tests
are performed over a shared random sample.

\section{Conclusion}
\label{sec:conclusion}

We studied the selection of $n$ tickets in a $(v,k)$-lottery to maximise
the probability that at least one ticket wins.  Three contributions stand
out.  First, for $n=2$ we proved that the joint prize probability is
non-increasing in pairwise overlap (Theorem~\ref{thm:two-ticket-monotone}),
so minimum-overlap collections are optimal.  For the Lotof\'{a}cil
specifically, any two tickets with overlap at most six have mutually exclusive
winning events and achieve $P_{\mathrm{win}} = 2p$ exactly
(Proposition~\ref{prop:mutual-excl}).  Second, two algorithms
(Algorithms~\ref{alg:greedy} and~\ref{alg:sa}) produce near-minimum-overlap
collections in time linear in $n$, with a provable $(1-1/e)$-approximation
guarantee for the greedy stage.  Third, a paired Monte Carlo experiment on
the Lotof\'{a}cil confirms gains of up to eleven percentage points over
random selection at $n=5$, with statistically significant advantages at
every $n \ge 2$ up to the saturation point near $n=50$.  These results
provide the first rigorous probabilistic treatment of multi-ticket lottery
selection in the high-$k$ regime and lay groundwork for extensions to
multi-tier prize structures and prize-division models.

\paragraph*{Acknowledgements.}  [To be completed.]

\appendix
\section{Computation of Table~\ref{tab:delta-g}}
\label{sec:appendix-computation}

The values $g(m)$ in Table~\ref{tab:delta-g} were computed by exact
enumeration of the multivariate hypergeometric distribution.  For parameters
$v=25$, $k=15$, $t=11$ and overlap $m$, the partition sizes are
$(|S_0|,|S_1|,|S_2|,|S_3|) = (m, 15{-}m, 15{-}m, m{-}5)$.  The joint-win
probability is
\[
  g(m) = \frac{1}{\binom{25}{15}}
  \sum_{\substack{x_0+x_1+x_2+x_3=15 \\
        x_i \le |S_i| \\
        x_0+x_1 \ge 11,\; x_0+x_2 \ge 11}}
  \binom{m}{x_0}\binom{15{-}m}{x_1}\binom{15{-}m}{x_2}\binom{m{-}5}{x_3}.
\]
All arithmetic used Python's \texttt{math.comb} for exact integer
evaluation; no floating-point rounding occurs until the final division
by $\binom{25}{15} = 3{,}268{,}760$.
